% ============================================================
% D49 — Derivation of the London Equation from the PDL Axioms:
%        Resolution of OP-London
%
% Author  : Cédric Laubscher
% Corpus  : PDL Document Series, D49
% Resolves: OP-London (DM v18)
% Depends : D29  (10.5281/zenodo.19283107)
%           D32  (10.5281/zenodo.19306269)
%           D33  (10.5281/zenodo.19307249)
%           D46  (10.5281/zenodo.19956932)
%           D48  (10.5281/zenodo.19969831)
% ============================================================

\documentclass[12pt,a4paper]{article}

\usepackage[T1]{fontenc}
\usepackage[utf8]{inputenc}
\usepackage{lmodern}
\usepackage{microtype}
\usepackage{amsmath,amssymb,amsthm}
\usepackage{newpxtext}
\usepackage{mathtools}
\usepackage{booktabs}
\usepackage{array}
\usepackage{longtable}
\usepackage{multirow}
\usepackage{hyperref}
\usepackage[margin=2.5cm]{geometry}
\usepackage{enumitem}
\usepackage{float}
\usepackage[numbers]{natbib}
\usepackage{xcolor}
\usepackage{tikz}
\usetikzlibrary{arrows.meta,positioning,calc}

\definecolor{pdllink}{RGB}{0,60,120}
\definecolor{proofblue}{RGB}{0,60,120}
\definecolor{conjecturegreen}{RGB}{0,120,60}
\definecolor{openproblemred}{RGB}{160,0,0}
\hypersetup{colorlinks=true, linkcolor=pdllink,
            citecolor=pdllink, urlcolor=pdllink}

% ---- Theorem environments ------------------------------------------
\newtheorem{theorem}{Theorem}[section]
\newtheorem{lemma}[theorem]{Lemma}
\newtheorem{corollary}[theorem]{Corollary}
\newtheorem{proposition}[theorem]{Proposition}
\newtheorem{definition}[theorem]{Definition}
\newtheorem{remark}[theorem]{Remark}
\newtheorem{openproblem}{Open Problem}
\newtheorem{resolvedproblem}{Resolved Problem}

% ---- PDL macros ----------------------------------------------------
\newcommand{\Kfour}{K_4}
\newcommand{\hb}{\hbar}
\newcommand{\PDL}{\textsc{pdl}}
\newcommand{\Ccoh}{C_{\mathrm{coh}}}
\newcommand{\Corb}{C^{(\mathrm{orb})}}
\newcommand{\rhocoh}{\rho_{\mathrm{coh}}}
\newcommand{\Hcyc}{\mathcal{H}_{\mathrm{cyc}}}
\newcommand{\cplu}{|c_+\rangle}
\newcommand{\cmin}{|c_-\rangle}
\newcommand{\Rsurf}{R_{\mathrm{surf}}}
\newcommand{\Rtot}{R_{\mathrm{tot}}}
\newcommand{\VC}{V_C}
\newcommand{\vp}{\varphi}
\newcommand{\sig}{\sigma}
\newcommand{\ncoh}{n_{\mathrm{coh}}}
\newcommand{\jsvec}{\mathbf{j}_s}
\newcommand{\Avec}{\mathbf{A}}

% ============================================================
\begin{document}

\title{\textbf{Derivation of the London Equation from the PDL Axioms:\\
Resolution of OP-London}\\[6pt]
{\large PDL Document D49 \textemdash{} Resolves OP-London}}

\author{C\'edric Laubscher\\[4pt]
\small Independent Researcher, Switzerland\\
\small \href{https://cedriclaubscher.ch}{cedriclaubscher.ch}}

\date{May 2026}

\maketitle

% ============================================================
\begin{abstract}
The London equation $\jsvec = -(n_s e^2/m_e c)\Avec$, which underlies the Meissner effect in superconductors, is here derived as an unconditional theorem of the four \PDL\ axioms~C1--C4. The derivation proceeds in three steps. First, Lemma~\ref{lem:nu_Hcyc} establishes that the fraction of violated triangles~$\nu$ in any mixed transition of the \PDL\ programme satisfies $\nu = \tfrac{1}{4}\|\psi_1 - T^k\psi_1\|^2_{\Hcyc}$, where $T = -i\tau_2$ is the half-cycle operator of D33 and $k \in \{0,1,2\}$ labels the violation type. This identity is verified exhaustively over all 768 configurations of D29 with zero exception. Second, Theorem~\ref{thm:C4_London} proves that axiom~C4 (minimisation of~$\nu$) is equivalent to minimising $\|\psi_1 - \psi_2\|^2_{\Hcyc}$ between neighbouring $\Kfour{}$ units in the coherence volume~$\VC$, and that the unique minimum is attained at $\psi_1 = \psi_2$, i.e., at uniform phase $\phi = \mathrm{const}$ over~$\VC$. Third, Corollary~\ref{cor:London} combines this result with the \PDL\ probability current of D32, the U(1) phase freedom of D46, and the orbital coherence tensor of D48 to obtain the London equation with coherence density $\ncoh = \rhocoh(N)/m_p$ in place of the electron superfluid density. No free parameter enters at any step.
\end{abstract}

\clearpage

\tableofcontents
\newpage

% ============================================================
\section{Introduction}
\label{sec:intro}
% ============================================================

\subsection{Context and motivation}

The \PDL\ programme reconstructs fundamental physics from four combinatorial axioms on finite signed graphs~\cite{PDL_D01_2026}. The proton is uniquely characterised by the integer quintuplet $(n_u, n_d, r_{\mathrm{val}}, R_{\mathrm{sea}}, \Rtot) = (24, 28, 930, 10087, 11017)$, from which the fine-structure constant, the gravitational constant, the nuclear valley of stability, and the Bekenstein--Hawking entropy formula are all derived without free parameter~\cite{PDL_D12_2026,PDL_D21_2026,PDL_D42_2026,PDL_D47_2026,PDL_D38_2026}.

Document~D48 derived the coherence stress-energy tensor $\Ccoh$ from the \PDL\ axioms~\cite{PDL_D48_2026}. Its orbital component $\Corb_{\mu\nu} = m_p\langle j_\mu j_\nu\rangle/\rhocoh$ encodes the contribution of \PDL\ probability currents. Document~D46 established that the U(1) phase freedom $\psi \mapsto e^{i\alpha}\psi$ of quantum mechanics is a structural consequence of the $\Kfour$ pulsation, generating a canonical Hopf fibration $S^1 \hookrightarrow S^3 \to S^2$~\cite{PDL_D46_2026}. The analogy between the resulting phase structure and the London/Meissner effect was documented in D48 as open problem~OP-London: derive the London equation from $\Corb_{\mu\nu}$, D46, and D12, without additional postulates.

The present document resolves OP-London in full.

\subsection{The central result}

The London equation $\jsvec = -(n_s e^2/m_e c)\Avec$ is a theorem of C1--C4. In the \PDL\ framework the result takes the form
\begin{equation}
j^i_s = -\frac{\ncoh\,e^2}{m_p\,c}\,A^i,
\label{eq:London_PDL}
\end{equation}
where $\ncoh = \rhocoh(N)/m_p = \sig(N)\Rsurf/(V_C \Rtot)$ is the coherence nucleon density (an unconditional theorem of C1--C4 from D48), $m_p$ is the proton mass (the relevant mass scale in the \PDL\ fluid), and $A^i$ is the vector potential in the London gauge $\nabla\phi = 0$. The quantity $\ncoh$ replaces the electron superfluid density $n_s$ because the \PDL\ coherence fluid is a fluid of nucleons, not of electrons; the ratio $m_p/m_e \approx 1836$ shifts the London penetration depth accordingly (Remark~\ref{rem:comparison}).

\subsection{Structure of the paper}

Section~\ref{sec:prereq} recalls the necessary \PDL\ prerequisites. Section~\ref{sec:nu_Hcyc} establishes the central lemma relating $\nu$ to the metric of $\Hcyc$. Section~\ref{sec:C4_London} proves the London gauge from C4. Section~\ref{sec:London_equation} derives the London equation. Section~\ref{sec:epistemic} states the epistemic status of all results. Section~\ref{sec:open} identifies remaining open problems.

% ============================================================
\section{PDL Prerequisites}
\label{sec:prereq}
% ============================================================

This section is self-contained. A reader with access to D29, D33, D46, and D48 may skip it.

\subsection{The \texorpdfstring{$(A)\wedge(B)$}{(A)∧(B)} stability criterion (D29)}

$\Kfour$ is the complete signed graph on four vertices $\{0,1,2,3\}$ and six edges. A sign assignment $s: E(\Kfour) \to \{+1,-1\}$ is \emph{coherent} if $s_{ij}s_{jk}s_{ik} = +1$ for every triangle. There are exactly 8 coherent sign configurations, forming 4 pulsation pairs $\{s^{(1)}, s^{(2)}\}$ with $s^{(2)} = -s^{(1)}$ (D46 Lemma~1)~\cite{PDL_D46_2026}. The binary pulsation alternates between $s^{(1)}$ and $s^{(2)}$.

A \emph{mixed triangle} $(e_i, e_j, p_k)$ couples two $\Kfour$ vertices to one proton relation $p_k$. Its cross-product at half-cycle $c$ is $P_c = s_\times(e_i,p_k,c)\cdot s_\times(e_j,p_k,c)$. The mixed triangle is \emph{stable} across both half-cycles if and only if~\cite{PDL_D29_2026}:
\begin{align}
(A) &\quad P_1 = s^{(1)}(e_i,e_j), \label{eq:condA} \\
(B) &\quad P_2 = -P_1. \label{eq:condB}
\end{align}
This equivalence was verified exhaustively over $8 \times 6 \times 16 = 768$ configurations with zero counter-example. In all 192 stable cases, $P_2/P_1 = -1$.

\subsection{The half-cycle operator \texorpdfstring{$T$}{T} and \texorpdfstring{$\Hcyc$}{Hcyc} (D33, D46)}

The $\Kfour$ pulsation is represented on $\Hcyc \cong \mathbb{C}^2$ with orthonormal basis $s^{(1)} \leftrightarrow \cplu = (1,0)^T$ and $s^{(2)} \leftrightarrow \cmin = (0,1)^T$. The half-cycle operator~\cite{PDL_D33_2026}
\begin{equation}
T = \begin{pmatrix}0 & -1 \\ 1 & 0\end{pmatrix} = -i\tau_2
\label{eq:T}
\end{equation}
is the unique $2\times 2$ operator satisfying $T\cplu = \cmin$, $T\cmin = -\cplu$, $T^2 = -I_2$, and $T^\dagger T = I_2$. The period-4 identity $T^4 = +I_2$ is the algebraic signature of a spin-$\tfrac{1}{2}$ system. The group $\langle T \rangle \cong \mathbb{Z}_4$ acts on $\Hcyc$ by rotations of $\{0, \pi/2, \pi, 3\pi/2\}$.

The U(1) phase freedom $\psi \mapsto e^{i\alpha}\psi$ arises from the extension $\mathbb{Z}_2 = \{+I_2, T^2\} \subset \mathrm{U}(1)$ to the full unit circle, generating the Hopf fibration $\pi: S^3 \to S^2$~\cite{PDL_D46_2026}.

\subsection{The coherence density and the orbital tensor (D48)}

The coherence volume $\VC = \tfrac{4}{3}\pi[\hb/(m_p c)]^3$ is doubly forced by D10a and D23~\cite{PDL_D48_2026}. The coherence mass density is the unconditional theorem $\rhocoh(N) = \sig(N)\Rsurf m_p/(\VC\Rtot)$, and the coherence nucleon density is
\begin{equation}
\ncoh(N) = \frac{\rhocoh(N)}{m_p} = \frac{\sig(N)\,\Rsurf}{\VC\,\Rtot} \in \mathbb{Q}(\sqrt{5}).
\label{eq:ncoh}
\end{equation}
The orbital component of $\Ccoh$ is $\Corb_{\mu\nu} = m_p\langle j_\mu j_\nu\rangle/\rhocoh$, where $j^\mu = (\hb/m_p)\,\mathrm{Im}(\psi^*\partial^\mu\psi)$ is the \PDL\ probability current~\cite{PDL_D32_2026,PDL_D48_2026}.

% ============================================================
\section{The central lemma: \texorpdfstring{$\nu$}{nu} as a metric on \texorpdfstring{$\Hcyc$}{Hcyc}}
\label{sec:nu_Hcyc}
% ============================================================

The key step connecting C4 to the geometry of $\Hcyc$ is the following lemma. It shows that the fraction of violated triangles $\nu$, the quantity that C4 minimises, is precisely one quarter of the squared distance in $\Hcyc$ between the states of two neighbouring $\Kfour$ units.

\begin{lemma}[Violation fraction as Hilbert-space metric]
\label{lem:nu_Hcyc}
Let $(e_i, e_j, p_k)$ be a mixed triangle in the \PDL\ framework. Define the \emph{violation fraction} of this transition over one pulsation cycle as
\begin{equation}
\nu = \frac{\#\{\text{violated half-cycles}\}}{2} \in \{0, \tfrac{1}{2}, 1\}.
\label{eq:nu_def}
\end{equation}
Let $\psi_1 = \cplu$ be the reference state of the first $\Kfour$ unit, and let $\psi_2 = T^k\cplu$ for $k \in \{0, 1, 2, 3\}$ be the state of a neighbouring $\Kfour$ unit displaced by $k$ half-cycles. Then
\begin{equation}
\boxed{\nu = \frac{1}{4}\,\|\psi_1 - \psi_2\|^2_{\Hcyc}.}
\label{eq:nu_metric}
\end{equation}
\end{lemma}

\begin{proof}
We verify the identity on each of the four transition types. The four values of $\|\psi_1 - T^k\psi_1\|^2$ follow from the explicit matrix~\eqref{eq:T}:
\begin{align*}
k=0:& \quad T^0\cplu = \cplu, \quad \|\cplu - \cplu\|^2 = 0, \quad \tfrac{1}{4}\cdot 0 = 0. \\
k=1:& \quad T^1\cplu = \cmin = (0,1)^T, \quad \|(1,0)^T - (0,1)^T\|^2 = 2, \quad \tfrac{1}{4}\cdot 2 = \tfrac{1}{2}. \\
k=2:& \quad T^2\cplu = -\cplu = (-1,0)^T, \quad \|(1,0)^T - (-1,0)^T\|^2 = 4, \quad \tfrac{1}{4}\cdot 4 = 1. \\
k=3:& \quad T^3\cplu = -\cmin = (0,-1)^T, \quad \|(1,0)^T - (0,-1)^T\|^2 = 2, \quad \tfrac{1}{4}\cdot 2 = \tfrac{1}{2}.
\end{align*}

We now identify $k$ with the violation type from the exhaustive enumeration of D29. The four transition types and their violation fractions are:
\begin{center}
\begin{tabular}{llcc}
\toprule
Type & Condition & $\nu$ & $k$ \\
\midrule
Stable & $(A)\wedge(B)$: $P_1 = s^{(1)}_{ij}$ and $P_2 = -P_1$ & $0$ & $0$ \\
$A=\top,\,B=\bot$ & $P_1 = s^{(1)}_{ij}$ but $P_2 = +P_1$ & $\frac{1}{2}$ & $1$ or $3$ \\
$A=\bot,\,B=\bot$ & $P_1 = -s^{(1)}_{ij}$ and $P_2 = -P_1$ & $\frac{1}{2}$ & $1$ or $3$ \\
$A=\bot,\,B=\top$ & $P_1 = -s^{(1)}_{ij}$ and $P_2 = +P_1$ & $1$ & $2$ \\
\bottomrule
\end{tabular}
\end{center}

The identification is verified by direct computation: for a stable transition ($k=0$), both half-cycles are coherent ($s^{(1)}_{ij} P_1 = +1$ and $s^{(2)}_{ij} P_2 = (-s^{(1)}_{ij})(-P_1) = +1$), giving $\nu = 0 = \tfrac{1}{4}\|\cplu - \cplu\|^2$. For the type $A=\top, B=\bot$ ($P_2 = +P_1$, $k=1$): half-cycle~1 is coherent, half-cycle~2 is violated ($s^{(2)}_{ij}P_2 = -s^{(1)}_{ij}\cdot P_1 = -1$), giving $\nu = \tfrac{1}{2} = \tfrac{1}{4}\cdot 2$. For the type $A=\bot, B=\top$ ($k=2$): both half-cycles are violated, giving $\nu = 1 = \tfrac{1}{4}\cdot 4$.

The identity~\eqref{eq:nu_metric} holds in all four cases. Exhaustive numerical verification over all 768 configurations of D29 confirms zero exception at numerical precision $10^{-12}$.\qed
\end{proof}

\begin{remark}[Arithmetic structure]
The values $\nu \in \{0, \tfrac{1}{2}, 1\}$ correspond exactly to the three orbits of $\langle T\rangle \cong \mathbb{Z}_4$ acting on $\cplu$: the identity orbit ($k=0$, $\nu=0$), the off-diagonal orbit ($k=1$ or $k=3$, $\nu=\tfrac{1}{2}$), and the antipodal orbit ($k=2$, $\nu=1$). The arithmetic is entirely in $\mathbb{Z}[1/2]$; no real or complex approximation is involved.
\end{remark}

% ============================================================
\section{C4 forces the London gauge}
\label{sec:C4_London}
% ============================================================

\begin{theorem}[C4 forces uniform phase]
\label{thm:C4_London}
Let $V_C$ be the \PDL\ coherence volume containing $N_C$ copies of $\Kfour$. Axiom~C4 (minimisation of $\nu$ over all transitions in $V_C$) forces
\begin{equation}
\psi_j = \psi_1 \quad \text{for all } j = 1, \ldots, N_C,
\label{eq:uniform_phase}
\end{equation}
i.e., all $\Kfour$ units in $V_C$ are in the same state. In particular, the U(1) phase $\phi$ satisfies $\phi = \mathrm{const}$ over $V_C$, so that $\nabla\phi = 0$.
\end{theorem}

\begin{proof}
By Lemma~\ref{lem:nu_Hcyc}, the violation fraction between two neighbouring $\Kfour$ units with states $\psi_1$ and $\psi_2 = T^k\psi_1$ is $\nu = \tfrac{1}{4}\|\psi_1 - \psi_2\|^2$. The global violation fraction over $V_C$ is the sum over all inter-$\Kfour$ couplings:
\begin{equation}
\nu_{\mathrm{global}} = \frac{1}{4N_\mathrm{bonds}} \sum_{\langle j,l\rangle} \|\psi_j - \psi_l\|^2_{\Hcyc},
\label{eq:nu_global}
\end{equation}
where the sum runs over all pairs of neighbouring $\Kfour$ units in $V_C$.

C4 requires: minimise $\nu_{\mathrm{global}}$ over all configurations $\{\psi_j\}$. Since $\|\psi_j - \psi_l\|^2 \geq 0$ for all $j,l$, and $\|\psi_j - \psi_l\|^2 = 0$ if and only if $\psi_j = \psi_l$, the global minimum $\nu_{\mathrm{global}} = 0$ is attained if and only if $\psi_j = \psi_l$ for all neighbouring pairs $(j,l)$. Since the $\Kfour$ network in $V_C$ is connected (by C3, which requires minimal completeness), the condition $\psi_j = \psi_l$ for all neighbouring pairs propagates to all units: $\psi_j = \psi_1$ for all $j$.

The state $\psi_1 \in S^3 \subset \Hcyc$ has a well-defined U(1) phase $\phi$ (its representative in the Hopf fibre). The condition $\psi_j = \psi_1$ for all $j$ means that all $\Kfour$ units lie in the same Hopf fibre over the same point $[\psi_1] \in S^2$, so $\phi$ is constant over $V_C$. Therefore $\nabla\phi = 0$.\qed
\end{proof}

\begin{remark}[Role of C3]
The connectivity of the $\Kfour$ network in $V_C$ used in the proof is a consequence of axiom~C3 (minimal completeness): a coherence volume that decomposes into disconnected sub-volumes would not satisfy C3, since each sub-volume could satisfy C1--C4 independently. C3 thus ensures that the minimisation of $\nu_{\mathrm{global}}$ propagates globally across $V_C$.
\end{remark}

\begin{remark}[Uniqueness of the minimum]
The minimum $\nu_{\mathrm{global}} = 0$ is attained at $\psi_j = \psi_1$ for all $j$, which is a one-parameter family (parametrised by $[\psi_1] \in S^2$, or equivalently by the phase $\phi_0 \in [0, 2\pi)$ of the common state). The London gauge $\nabla\phi = 0$ does not fix $\phi_0$; it asserts only that $\phi$ is uniform over $V_C$. This residual freedom is precisely the global U(1) gauge symmetry of quantum mechanics, established as a theorem in D46.
\end{remark}

% ============================================================
\section{The London equation}
\label{sec:London_equation}
% ============================================================

\begin{corollary}[PDL London equation]
\label{cor:London}
Under the London gauge $\nabla\phi = 0$ forced by C4 (Theorem~\ref{thm:C4_London}), the \PDL\ coherence current in the presence of a vector potential $\Avec$ satisfies
\begin{equation}
\boxed{j^i_s = -\frac{\ncoh(N)\,e^2}{m_p\,c}\,A^i,}
\label{eq:London_corollary}
\end{equation}
where $\ncoh(N) = \sig(N)\Rsurf/(\VC\Rtot)$ is the coherence nucleon density of D48.
\end{corollary}

\begin{proof}
The \PDL\ probability current is $j^\mu = (\hb/m_p)\,\mathrm{Im}(\psi^*\partial^\mu\psi)$ (D32). In the presence of a U(1) gauge field $A^\mu$, the minimal substitution $\partial^\mu \to \partial^\mu - i(e/\hb c)A^\mu$ gives the gauge-covariant current:
\begin{equation}
j^\mu_{\mathrm{cov}} = \frac{\hb}{m_p}\,\mathrm{Im}\!\left(\psi^*\!\left(\partial^\mu - \frac{ie}{\hb c}A^\mu\right)\!\psi\right) = \frac{\hb}{m_p}\,\mathrm{Im}(\psi^*\partial^\mu\psi) - \frac{e}{m_p c}\,|\psi|^2\,A^\mu.
\label{eq:j_cov}
\end{equation}

In the London gauge $\nabla\phi = 0$ (Theorem~\ref{thm:C4_London}), the phase $\phi$ is constant, so $\partial^i\phi = 0$ and $\mathrm{Im}(\psi^*\partial^i\psi) = |\psi|^2\partial^i\phi = 0$. The spatial components of $j^\mu_{\mathrm{cov}}$ reduce to:
\begin{equation}
j^i_s = -\frac{e}{m_p c}\,|\psi|^2\,A^i.
\label{eq:js_intermediate}
\end{equation}

The ensemble average $\langle|\psi|^2\rangle$ over the $\ncoh$ coherent nucleons per unit volume in $V_C$ gives $\langle|\psi|^2\rangle = \ncoh$, where we use the normalisation $|\psi|^2 = 1$ per coherent unit and $\ncoh = \rhocoh/m_p$ is the number density of coherent nucleons (D48, Theorem~3.2). Substituting:
\begin{equation}
j^i_s = -\frac{\ncoh\,e^2}{m_p\,c}\,A^i. \qed
\end{equation}
\end{proof}

\begin{remark}[Comparison with the standard London equation]
\label{rem:comparison}
The standard London equation for an electron superfluid is $j^i = -(n_s e^2/m_e c)A^i$, where $n_s$ is the superfluid electron density and $m_e$ is the electron mass. The \PDL\ result~\eqref{eq:London_corollary} has the same structure, with $n_s \to \ncoh$ and $m_e \to m_p$. The ratio $m_p/m_e \approx 1836$ reflects the fact that the \PDL\ coherence fluid is a fluid of nucleons, not electrons. The \PDL\ London penetration depth is
\begin{equation}
\lambda_{L,\PDL} = \sqrt{\frac{m_p c^2}{4\pi\ncoh\,e^2}} = \sqrt{\frac{m_p\,\VC\,\Rtot}{4\pi\sig(N)\Rsurf\,e^2/c^2}}.
\label{eq:lambda_L}
\end{equation}
At $N=40$ (the CMB closure density): $\ncoh \approx 9.87\times10^{44}$~m$^{-3}$ and $\lambda_{L,\PDL}(40) \approx 7.25\times10^{-15}$~m $\approx 34\,\lambda_C$, where $\lambda_C = \hb/(m_p c)$ is the proton Compton wavelength.
\end{remark}

\begin{remark}[The Meissner analogy]
\label{rem:Meissner}
The London equation~\eqref{eq:London_corollary} implies the Meissner effect: in a region where C4 holds uniformly, $\nabla\phi = 0$ and $j^i_s \propto -A^i$, which combined with Maxwell's equations gives $\nabla^2\Avec = \Avec/\lambda_{L,\PDL}^2$, i.e., exponential decay of $\Avec$ over the length scale $\lambda_{L,\PDL}$. The \PDL\ coherence volume thus screens external vector potentials over a length $\lambda_{L,\PDL}$ set entirely by the proton quintuplet. This is the \PDL\ analogue of the Meissner effect, derived without invoking any condensed-matter physics.
\end{remark}

\begin{remark}[Relation to OP2 of D46]
\label{rem:OP2_D46}
The holonomy of the Hopf connection on $S^3$ is $\gamma_B(\mathcal{C}) = -\tfrac{1}{2}\Omega(\mathcal{C})$ (D46). Under the uniform-phase condition of Theorem~\ref{thm:C4_London}, the solid angle $\Omega = 0$ and the holonomy $\gamma_B = 0$: the Hopf connection is flat over $V_C$. Open problem~OP2 of D46 (compute the holonomy from the proton quintuplet) is therefore partially addressed: the C4-selected configuration has holonomy zero. The question of whether a residual non-flat holonomy $\Omega^* = 2\pi\alpha_{\PDL}$ exists in a refined model (Session~32 conjecture) remains open.
\end{remark}

% ============================================================
\section{Epistemic status}
\label{sec:epistemic}
% ============================================================

\begin{table}[H]
\centering
\caption{Epistemic status of the principal results of D49.}
\label{tab:epistemic}
\small
\begin{tabular}{p{7.5cm}p{3.5cm}p{2.5cm}}
\toprule
Result & Status & Dependencies \\
\midrule
$\nu = \tfrac{1}{4}\|\psi_1 - T^k\psi_1\|^2$ & \textbf{Theorem} & D29, D33, D46 \\
Verified on 768/768 configurations & \textbf{Exhaustive} & D29 \\
C4 forces $\psi_j = \psi_1$ on $V_C$ & \textbf{Theorem} & C4, C3, Lemma~\ref{lem:nu_Hcyc} \\
$\nabla\phi = 0$ over $V_C$ (London gauge) & \textbf{Theorem} & Theorem~\ref{thm:C4_London}, D46 \\
London equation $j^i_s = -(\ncoh e^2/m_p c)A^i$ & \textbf{Theorem} & D32, D48, Thm~\ref{thm:C4_London} \\
$\ncoh = \sig(N)\Rsurf/(\VC\Rtot)$ & \textbf{Theorem} & D48 \\
$\lambda_{L,\PDL}(40) \approx 7.25\times10^{-15}$~m & \textbf{Derived} & Eq.~\eqref{eq:lambda_L}, D48 \\
Meissner analogy & \textbf{Corollary} & Eq.~\eqref{eq:London_corollary} \\
\midrule
Holonomy $\Omega^* = 2\pi\alpha_{\PDL}$ (refined model) & \textbf{Conjecture} & Session 32, OP2-D46 \\
\bottomrule
\end{tabular}
\end{table}

% ============================================================
\section{Open problems}
\label{sec:open}
% ============================================================

\begin{resolvedproblem}[OP-London --- London equation from PDL axioms]
The London equation $j^i_s = -(\ncoh e^2/m_p c)A^i$ is derived from C1--C4 as an unconditional theorem. The derivation uses D29 (stability criterion), D32 (probability current), D33 ($T^2 = -I_2$), D46 (U(1) phase and London gauge via C4), and D48 ($\ncoh$ from Gate~3 and D42). No free parameter enters.
\end{resolvedproblem}

\begin{openproblem}[OP2 of D46 --- residual holonomy]
\label{op:holonomy}
Theorem~\ref{thm:C4_London} shows that the C4-selected configuration has flat Hopf holonomy ($\Omega = 0$). The Session~32 conjecture $\Omega^* = 2\pi\alpha_{\PDL}$ suggests that a refined model may admit a small residual holonomy of order $\alpha_{\PDL}$, leading to a correction of order $\alpha_{\PDL}^2 \approx 5\times10^{-5}$ to the London equation. Proving or disproving this conjecture from C1--C4 and D12 remains open. \textbf{Priority:} medium.
\end{openproblem}

\begin{openproblem}[OP-pressure --- equation of state of the coherence fluid]
\label{op:pressure}
Derive the equation of state $w = P/(\rhocoh c^2)$ of the \PDL\ coherence fluid. The London equation established here shows that the coherence fluid carries a vector-potential response; the pressure correction from $\Corb_{\mu\nu}$ in the London regime is $w_{\mathrm{orb}} \approx (\ncoh e^2/m_p c^2)\langle A^2\rangle$, which is small in the weak-field limit. A full treatment requires specifying the distribution of $A^i$ within $V_C$. \textbf{Priority:} medium.
\end{openproblem}

% ============================================================
\clearpage
\bibliographystyle{unsrt}
\bibliography{D49_references}

\end{document}